\documentclass[11pt,a4paper]{article}
\usepackage[utf8]{inputenc}
\usepackage{amsmath,amssymb,amsthm}
\usepackage{hyperref}
\usepackage{booktabs}
\usepackage{cleveref}

% === Fill in these fields ===
\title{[Title: Bottom-Line Result in One Phrase]}
\author{Ramanujan Agent \\
  \textit{Ramanujan Machine Project}}
\date{\today}

% === Notation macros (extend as needed) ===
\newcommand{\pcf}{\mathrm{PCF}}
\newcommand{\cmf}{\mathrm{CMF}}

\begin{document}
\maketitle

% ============================================================
\section*{Bottom Line}
% ============================================================
% One paragraph. State the main result. Be specific and quantitative.
% Example: "We show that the CMF trajectory [0,1,1,1]/[1,1,1] applied to
%  the Ramanujan-type 4F3 family yields a new continued fraction for 1/pi
%  converging at 3.01 digits per term, verified to 500 decimal places."

\textbf{TODO:} State the main result here.

% ============================================================
\section{Setup and Notation}
% ============================================================
% Define all symbols used. Keep it short — use a table if many symbols.
% Reference the CMF/PCF framework from the Ramanujan Machine papers.

% ============================================================
\section{Main Result}
% ============================================================
% Formal statement of the formula, identity, theorem, or computation.
% Include the full mathematical expression.

% ============================================================
\section{Worked Examples}
% ============================================================
% At least 2 numerical examples showing the result in action.
% Include intermediate computation steps.

\subsection{Example 1}
% State the parameters, compute, show the result matches the expected constant.

\subsection{Example 2}
% Different parameters or a different aspect of the result.

% ============================================================
\section{Numerical Verification}
% ============================================================
% Describe the verification procedure:
% - What code was run (include filename or inline snippet)
% - How many digits were verified
% - What reference values were used
% - Any edge cases tested
%
% Example:
% "Computed at depth N=1000 using mpmath with dps=2000.
%  The convergent matches mpmath.pi to 503 decimal places."

% ============================================================
\section{Context and Classification}
% ============================================================
% How does this relate to known CMFs, PCFs, or prior results?
% Is it coboundary-equivalent to something known?
% What is the irrationality measure delta?

% ============================================================
\section{Open Questions}
% ============================================================
% What remains unproven? What would strengthen the result?
% Be honest about limitations.

% ============================================================
% \section*{Appendix: Verification Code}
% ============================================================
% (Optional) Include the Python code used for verification.
% \begin{verbatim}
% import mpmath
% mpmath.mp.dps = 500
% ...
% \end{verbatim}

\end{document}
